(1) 转动自由度配分是 $\boxed{Z_r(\beta)=\sum_{J=0}^{\infty}(2J+1)e^{-\beta k_B\Theta_rJ(J+1)}}$，其中 $\beta=1/(k_BT)$。$2J+1$ 是空间角动量磁量子数的简并，而非人为额外的可占位置数。

(2) 平均能量 $U_r=-\partial_\beta\ln Z_r$，定容热容 $C_r=k_B\beta^2(\langle E_J^2\rangle-\langle E_J\rangle^2)\ge0$。低温只须保留 $J=0$ 单态及 $J=1$ 三重态：设 $x=\Theta_r/T\gg1$，则 $Z_r\simeq1+3e^{-2x}$，$U_r\simeq6k_B\Theta_re^{-2x}$，$\boxed{C_r\simeq12k_Bx^2e^{-2x}}$。

(3) 高温时把求和近似为积分：$\int_0^\infty(2J+1)e^{-xJ(J+1)}dJ=1/x=T/\Theta_r$。因此 $\ln Z_r\simeq\ln(T/\Theta_r)$、$U_r\simeq k_BT$、$\boxed{C_r\to k_B}$，对应绕分子轴以外的两个转动二次型各贡献 $k_BT/2$。

(4) 真正同核分子会受到核自旋和总波函数交换对称性的约束，使奇、偶 $J$ 态带不同统计权重，低温第一允许激发能隙与简并可能改变。只要基态有限简并且最低激发有正能隙，非残余熵的转动热容趋零，与第三定律一致。
